% ======================================================================
% CHAPTER 5: DEEP REINFORCEMENT LEARNING AGENT ENGINEERING
% ======================================================================
\chapter{Deep Reinforcement Learning Agent Engineering}
\label{chap:agent}

\section{Dueling Double Deep Q-Network Architecture}
The core intelligence of our traffic engineering platform resides in the \texttt{DQNRoutingAgent} class (\texttt{agent/dqn\_router.py}), which implements a PyTorch \textbf{Dueling Double Deep Q-Network (D3QN)} stabilized by \textbf{Layer Normalization} and an $O(\log N)$ \textbf{Binary SumTree Prioritized Experience Replay} engine.

\subsection{Architectural Specification and Layer Dimensions}
The neural network takes the normalized 10-dimensional state vector $s_t \in \mathbb{R}^{10}$ as input and outputs estimated action values $Q(s_t, a_t)$ for each of the $K=4$ candidate paths. The network is structured into three distinct computational stages: a shared feature extraction trunk, a scalar state-value stream $V(s)$, and a vector action-advantage stream $A(s, a)$.

\begin{figure}[htbp]
\centering
\begin{minipage}{0.92\textwidth}
\centering
\fbox{\parbox{0.95\textwidth}{
\vspace{0.2cm}
\centering
{\large \textbf{PyTorch Dueling Q-Network Layer Topology}}\\[0.3cm]
\begin{tabular}{c}
\fbox{\parbox{0.75\textwidth}{\centering \textbf{Input State Vector} $s_t \in \mathbb{R}^{10}$\\ \footnotesize (Normalized hop count, 4 candidate bottleneck loads, 4 path delays, global max load)}}\\[0.2cm]
$\Downarrow$\\[0.2cm]
\fbox{\parbox{0.75\textwidth}{\centering \textbf{Shared Feature Extraction Trunk}\\ \footnotesize $\text{Linear}(10 \to 64) \to \text{LayerNorm}(64) \to \text{ReLU}$\\ $\text{Linear}(64 \to 128) \to \text{LayerNorm}(128) \to \text{ReLU}$}}\\[0.2cm]
$\swarrow$ \hspace{4cm} $\searrow$\\[0.2cm]
\begin{tabular}{c c}
\fbox{\parbox{0.40\textwidth}{\centering \textbf{Value Stream $V(s)$}\\ \footnotesize $\text{Linear}(128 \to 64)$\\ $\text{LayerNorm}(64) \to \text{ReLU}$\\ $\text{Linear}(64 \to 1) \implies V(s) \in \mathbb{R}^1$}} &
\fbox{\parbox{0.40\textwidth}{\centering \textbf{Advantage Stream $A(s, a)$}\\ \footnotesize $\text{Linear}(128 \to 64)$\\ $\text{LayerNorm}(64) \to \text{ReLU}$\\ $\text{Linear}(64 \to 4) \implies A(s, a) \in \mathbb{R}^4$}}
\end{tabular}\\[0.3cm]
$\searrow$ \hspace{4cm} $\swarrow$\\[0.2cm]
\fbox{\parbox{0.82\textwidth}{\centering \textbf{Identifiability Aggregation Layer}\\ \footnotesize $Q(s, a) = V(s) + \left( A(s, a) - \frac{1}{4}\sum_{a'=0}^3 A(s, a') \right) \implies Q(s, a) \in \mathbb{R}^4$}}\\
\end{tabular}
\vspace{0.2cm}
}}
\end{minipage}
\caption{Detailed Computational Architecture of the Dueling Q-Network with Layer Normalization.}
\label{fig:nn_detailed_topology}
\end{figure}

Table~\ref{tab:layer_specifications} details the parameter counts, weight tensor dimensions, and activation functions across each layer of the model.

\begin{table}[htbp]
\centering
\caption{Layer Configuration and Parameter Counts of the Dueling Q-Network}
\label{tab:layer_specifications}
\small
\begin{tabularx}{\textwidth}{l c c c X}
\toprule
\textbf{Layer Name} & \textbf{Input Dim} & \textbf{Output Dim} & \textbf{Parameters} & \textbf{Activation / Normalization} \\
\midrule
\texttt{shared\_fc1} & 10 & 64 & $10 \times 64 + 64 = 704$ & \texttt{nn.LayerNorm(64)} + ReLU \\
\texttt{shared\_fc2} & 64 & 128 & $64 \times 128 + 128 = 8,320$ & \texttt{nn.LayerNorm(128)} + ReLU \\
\addlinespace
\texttt{val\_fc1} & 128 & 64 & $128 \times 64 + 64 = 8,256$ & \texttt{nn.LayerNorm(64)} + ReLU \\
\texttt{val\_out} & 64 & 1 & $64 \times 1 + 1 = 65$ & Linear (Scalar $V(s)$) \\
\addlinespace
\texttt{adv\_fc1} & 128 & 64 & $128 \times 64 + 64 = 8,256$ & \texttt{nn.LayerNorm(64)} + ReLU \\
\texttt{adv\_out} & 64 & 4 & $64 \times 4 + 4 = 260$ & Linear (Vector $A(s, a)$) \\
\midrule
\multicolumn{3}{l}{\textbf{Total Trainable Parameters}} & \textbf{25,861} & ($\sim 101\text{ KB}$ memory footprint) \\
\bottomrule
\end{tabularx}
\end{table}

\subsection{The Role of Layer Normalization in Network Control Loops}
In computer network control environments, link loads, queuing delays, and packet dropouts exhibit extreme spatial and temporal volatility. An elephant flow arrival can cause a link's measured utilization to jump from $0.05$ to $0.98$ within a single polling interval. In standard unnormalized networks, such sharp input shifts cause internal covariate shift, resulting in exploding gradients that destabilize training or collapse Q-values into local minima.

To ensure numerical stability, we incorporate \textbf{Layer Normalization} (\texttt{nn.LayerNorm}) \cite{ba2016layer} after each dense linear projection. Unlike Batch Normalization (which computes statistics across mini-batch samples and fails during single-sample real-time inference), Layer Normalization standardizes activations across the channel dimension for each transition independently:
\begin{equation}
    \mu_L = \frac{1}{H} \sum_{i=1}^H x_i, \quad \sigma_L^2 = \frac{1}{H} \sum_{i=1}^H (x_i - \mu_L)^2
\end{equation}
\begin{equation}
    \text{LN}(x) = \frac{x - \mu_L}{\sqrt{\sigma_L^2 + \epsilon}} \odot \gamma + \beta
\end{equation}
where $H$ is the hidden layer dimension, and $\gamma, \beta \in \mathbb{R}^H$ are learnable affine parameters. Layer Normalization guarantees that activations remain strictly zero-centered and unit-variance, allowing smooth gradient descent even during violent traffic surges.

\subsection{Mathematical Proof of Dueling Identifiability}
A fundamental challenge in decomposing action values into value $V(s)$ and advantage $A(s, a)$ is \emph{unidentifiability}. Given $Q(s, a) = V(s) + A(s, a)$, adding an arbitrary constant $c$ to $V(s)$ while subtracting $c$ from $A(s, a)$ produces the exact same $Q$-value, preventing unique parameter recovery.

\begin{theorem}[Identifiability of Mean-Subtracted Dueling Network]
The aggregation formula:
\begin{equation}
    Q(s, a; \theta, \alpha, \beta) = V(s; \theta, \beta) + \left( A(s, a; \theta, \alpha) - \frac{1}{|\mathcal{A}|} \sum_{a' \in \mathcal{A}} A(s, a'; \theta, \alpha) \right)
    \label{eq:dueling_identifiability_proof}
\end{equation}
guarantees unique identifiability such that $V(s)$ represents the arithmetic mean of optimal action values across the action space.
\end{theorem}

\begin{proof}
Let $\bar{A}(s) = \frac{1}{|\mathcal{A}|} \sum_{a' \in \mathcal{A}} A(s, a')$. Summing Equation~\ref{eq:dueling_identifiability_proof} over all actions $a \in \mathcal{A}$:
\begin{align}
    \sum_{a \in \mathcal{A}} Q(s, a) &= \sum_{a \in \mathcal{A}} \left( V(s) + A(s, a) - \bar{A}(s) \right) \nonumber \\
    &= |\mathcal{A}| \cdot V(s) + \sum_{a \in \mathcal{A}} A(s, a) - |\mathcal{A}| \cdot \bar{A}(s) \nonumber \\
    &= |\mathcal{A}| \cdot V(s) + |\mathcal{A}| \bar{A}(s) - |\mathcal{A}| \bar{A}(s) \nonumber \\
    &= |\mathcal{A}| \cdot V(s)
\end{align}
Dividing both sides by $|\mathcal{A}|$:
\begin{equation}
    V(s) = \frac{1}{|\mathcal{A}|} \sum_{a \in \mathcal{A}} Q(s, a)
\end{equation}
Thus, $V(s)$ is uniquely identified as the average value of all candidate actions in state $s$, and the advantage terms $A(s, a) - \bar{A}(s)$ uniquely represent zero-mean deviations, establishing mathematical identifiability.
\end{proof}

\section{Binary SumTree Prioritized Experience Replay (PER)}
To overcome the limitations of uniform replay sampling, the \texttt{PrioritizedReplayBuffer} (\texttt{agent/prioritized\_replay.py}) implements an iterative binary \textbf{SumTree} data structure supporting $O(\log N)$ sampling and priority updates.

\subsection{Array-Based Binary SumTree Data Structure}
A binary SumTree is a complete binary tree where:
\begin{itemize}
    \item Each leaf node stores the priority $p_i$ of an individual transition tuple $e_i = (s_i, a_i, r_i, s_{i+1})$.
    \item Each internal node stores the exact sum of the priorities of its two child subtrees:
    \begin{equation}
        \text{tree}[\text{parent}] = \text{tree}[\text{left\_child}] + \text{tree}[\text{right\_child}]
    \end{equation}
    \item The root node ($\text{tree}[0]$) stores the grand sum of all transition priorities in memory:
    \begin{equation}
        \text{tree}[0] = \sum_{k=1}^N p_k
    \end{equation}
\end{itemize}

For a buffer capacity of $N = 5,000$ transitions, the SumTree is allocated as a single, contiguous 1D NumPy array of size $2N - 1 = 9,999$ elements:
\begin{itemize}
    \item Internal parent nodes occupy indices $0$ through $N - 2$.
    \item Leaf nodes occupy indices $N - 1$ through $2N - 2$.
    \item For any internal node at index $i$:
    \begin{equation}
        \text{left\_child}(i) = 2i + 1, \quad \text{right\_child}(i) = 2i + 2
    \end{equation}
    \item For any child node at index $j > 0$:
    \begin{equation}
        \text{parent}(j) = \lfloor (j - 1) / 2 \rfloor
    \end{equation}
\end{itemize}

Figure~\ref{fig:sumtree_diagram} illustrates the structural memory mapping of the binary SumTree.

\begin{figure}[htbp]
\centering
\begin{minipage}{0.85\textwidth}
\centering
\fbox{\parbox{0.95\textwidth}{
\vspace{0.2cm}
\centering
{\large \textbf{Array-Based Binary SumTree Structure ($N=4$)}}\\[0.3cm]
\begin{tabular}{c c c c}
\multicolumn{4}{c}{\textbf{Root Node}: $\text{tree}[0] = 42.0$ (Total Priority Sum)}\\[0.2cm]
\multicolumn{2}{c}{$\swarrow$} & \multicolumn{2}{c}{$\searrow$}\\[0.1cm]
\multicolumn{2}{c}{\textbf{Node 1}: $\text{tree}[1] = 29.0$} & \multicolumn{2}{c}{\textbf{Node 2}: $\text{tree}[2] = 13.0$}\\[0.2cm]
$\swarrow$ & $\searrow$ & $\swarrow$ & $\searrow$\\[0.1cm]
\textbf{Leaf 3}: $p_1 = 12.0$ & \textbf{Leaf 4}: $p_2 = 17.0$ & \textbf{Leaf 5}: $p_3 = 3.0$ & \textbf{Leaf 6}: $p_4 = 10.0$\\
\end{tabular}\\[0.3cm]
\rule{0.8\textwidth}{0.4pt}\\[0.2cm]
\textbf{Contiguous NumPy Array Representation}:\\
\texttt{tree = [42.0, 29.0, 13.0, 12.0, 17.0, 3.0, 10.0]}\\
\vspace{0.1cm}
}}
\end{minipage}
\caption{Binary SumTree Data Structure with Total Prefix Sum at Root.}
\label{fig:sumtree_diagram}
\end{figure}

\subsection{Sampling and Update Algorithms in Formal Pseudocode}
Sampling a mini-batch of size $B=32$ transitions proceeds by partitioning the range $[0, \text{tree}[0]]$ into $B$ equal intervals:
\begin{equation}
    I_k = \left[ k \cdot \frac{\text{tree}[0]}{B}, \, (k + 1) \cdot \frac{\text{tree}[0]}{B} \right) \quad \forall k \in \{0, 1, \dots, B-1\}
\end{equation}
A uniform random value $v_k \sim \mathcal{U}(I_k)$ is drawn from each interval, and the tree is traversed from the root downward to locate the corresponding leaf in exactly $\lceil \log_2 N \rceil = 13$ steps. Algorithm~\ref{alg:sumtree_sampling} and Algorithm~\ref{alg:sumtree_update} detail the iterative, non-recursive implementations.

\begin{algorithm}[htbp]
\caption{Iterative Non-Recursive Binary SumTree Retrieval}
\label{alg:sumtree_sampling}
\begin{algorithmic}[1]
\REQUIRE SumTree array $\text{tree}$ of size $2N - 1$, search value $v \in [0, \text{tree}[0]]$
\ENSURE Leaf index $\text{tree\_idx}$, priority value $p$, and data pointer $\text{data\_idx}$
\STATE $\text{idx} \gets 0$ \COMMENT{Start at root node}
\WHILE{$\text{idx} < N - 1$} \COMMENT{Loop until reaching a leaf node}
    \STATE $\text{left} \gets 2 \cdot \text{idx} + 1$
    \STATE $\text{right} \gets \text{left} + 1$
    \IF{$v \le \text{tree}[\text{left}]$}
        \STATE $\text{idx} \gets \text{left}$
    \ELSE
        \STATE $v \gets v - \text{tree}[\text{left}]$
        \STATE $\text{idx} \gets \text{right}$
    \ENDIF
\ENDWHILE
\STATE $\text{data\_idx} \gets \text{idx} - (N - 1)$
\RETURN $\text{idx}, \, \text{tree}[\text{idx}], \, \text{data\_idx}$
\end{algorithmic}
\end{algorithm}

\begin{algorithm}[htbp]
\caption{Iterative Non-Recursive Binary SumTree Priority Update}
\label{alg:sumtree_update}
\begin{algorithmic}[1]
\REQUIRE SumTree array $\text{tree}$, leaf index $\text{tree\_idx}$, new priority $p_{\text{new}}$
\ENSURE Updated prefix sums propagated from leaf to root
\STATE $\Delta \gets p_{\text{new}} - \text{tree}[\text{tree\_idx}]$ \COMMENT{Compute priority differential}
\STATE $\text{tree}[\text{tree\_idx}] \gets p_{\text{new}}$
\STATE $\text{idx} \gets \text{tree\_idx}$
\WHILE{$\text{idx} > 0$} \COMMENT{Propagate differential up to root}
    \STATE $\text{idx} \gets \lfloor (\text{idx} - 1) / 2 \rfloor$ \COMMENT{Move to parent}
    \STATE $\text{tree}[\text{idx}] \gets \text{tree}[\text{idx}] + \Delta$
\ENDWHILE
\end{algorithmic}
\end{algorithm}

\subsection{Importance Sampling Weight Annealing}
Prioritized sampling introduces bias because transitions with high TD error are sampled more frequently than their true empirical distribution. We correct this bias using Importance Sampling (IS) weights:
\begin{equation}
    w_i = \left( \frac{1}{N} \cdot \frac{1}{P(i)} \right)^\beta = \left( \frac{\text{tree}[0]}{N \cdot p_i^\alpha} \right)^\beta
\end{equation}
Weights are normalized by the maximum weight in the mini-batch:
\begin{equation}
    \tilde{w}_i = \frac{w_i}{\max_j w_j} \in (0.0, 1.0]
\end{equation}
The parameter $\beta$ controls the degree of correction. Early in training, gradient direction is noisy, so partial correction ($\beta_0 = 0.4$) accelerates discovery. As policy convergence approaches, full correction ($\beta \to 1.0$) is essential for asymptotic unbiased convergence:
\begin{equation}
    \beta(t) = \min\left(1.0, \, \beta_0 + t \cdot \frac{1.0 - \beta_0}{T_{\text{max}}}\right)
\end{equation}

\section{Polyak Soft Target Updates vs. Periodic Hard Copies}
In standard DQN \cite{mnih2015human}, target network parameters $\theta^-$ are updated via periodic hard copies every $C$ steps ($\theta^- \gets \theta$). In network traffic engineering, hard updates trigger sudden, violent shifts in target values $Y_t$, causing transient routing oscillations and policy collapse immediately following each synchronization.

To eliminate target instability, our agent implements \textbf{Polyak Soft Target Updates} \cite{lillicrap2015continuous}. After every gradient backpropagation step, target parameters are blended continuously toward online parameters:
\begin{equation}
    \theta^- \gets \tau \cdot \theta + (1 - \tau) \cdot \theta^-
    \label{eq:polyak_update}
\end{equation}
where $\tau = 0.005$ is the soft tracking rate. Polyak averaging acts as an exponential low-pass filter over parameter trajectories, ensuring that target Q-values evolve smoothly and monotonically.

\section{Huber Loss vs. Mean Squared Error}
Standard Mean Squared Error (MSE), $\mathcal{L}_{\text{MSE}}(\delta) = \delta^2$, has a quadratic gradient with respect to TD error: $\nabla_\theta \mathcal{L} = 2 \delta \nabla_\theta Q$. When sudden traffic bursts occur, the instantaneous TD error $\delta_i$ can exceed $50.0$, generating massive gradients ($2 \times 50 = 100$) that destabilize the Adam optimizer.

We employ the \textbf{Smooth L1 (Huber) Loss Function}:
\begin{equation}
    \mathcal{L}_{\text{Huber}}(\delta) = \begin{cases}
        0.5 \cdot \delta^2 & \text{if } |\delta| \le 1.0 \\[0.5em]
        |\delta| - 0.5 & \text{otherwise}
    \end{cases}
\end{equation}
whose gradient is:
\begin{equation}
    \frac{\partial \mathcal{L}_{\text{Huber}}}{\partial \delta} = \begin{cases}
        \delta & \text{if } |\delta| \le 1.0 \\
        \text{sign}(\delta) & \text{otherwise}
    \end{cases}
\end{equation}
For errors exceeding $1.0$, the gradient magnitude is strictly capped at $\pm 1.0$. This linear clipping prevents exploding gradients, guaranteeing robust convergence under extreme network traffic volatility.

\section{Complete End-to-End D3QN-PER Training Algorithm}
Algorithm~\ref{alg:d3qn_complete_loop} formalizes the complete training loop executed across the 1,000 curriculum episodes.

\begin{algorithm}[htbp]
\caption{Complete D3QN-PER Adaptive Traffic Engineering Training Pipeline}
\label{alg:d3qn_complete_loop}
\begin{algorithmic}[1]
\REQUIRE Mininet network topology $G = (V, E)$, episodes $M = 1000$, steps per episode $T = 30$, batch size $B = 32$, replay capacity $N = 5000$, learning rate $\alpha = 0.001$, soft update rate $\tau = 0.005$, discount factor $\gamma = 0.95$
\ENSURE Optimized D3QN weights $\theta^*$
\STATE Initialize online network parameters $\theta$ randomly; target network $\theta^- \gets \theta$
\STATE Initialize SumTree replay buffer $\mathcal{D}$ with capacity $N$, priorities $\alpha = 0.6$, $\beta_0 = 0.4$
\FOR{episode $m = 1$ \TO $M$}
    \STATE Calculate curriculum exploration rate $\epsilon(m)$ and injection stress profile
    \STATE Reset network link states; observe initial state $s_0 \in \mathbb{R}^{10}$
    \FOR{step $t = 1$ \TO $T$}
        \STATE Sample $r \sim \mathcal{U}(0, 1)$
        \IF{$r < \epsilon(m)$}
            \STATE Select random candidate path action $a_t \sim \mathcal{U}(\{0, 1, 2, 3\})$
        \ELSE
            \STATE Select greedy action $a_t = \arg\max_{a} Q(s_t, a; \theta)$ via forward pass
        \ENDIF
        \STATE Route flow along path $\mathcal{P}_{a_t}$; accumulate link traffic
        \STATE Calculate multi-objective reward $r_t = \mathcal{R}(s_t, a_t)$ via Equation~\ref{eq:multi_objective_reward}
        \STATE Observe subsequent state $s_{t+1} \in \mathbb{R}^{10}$
        \STATE Store transition $(s_t, a_t, r_t, s_{t+1})$ in SumTree $\mathcal{D}$ with max priority $p_{\text{max}}$
        \IF{$|\mathcal{D}| \ge B$}
            \STATE Sample mini-batch of $B$ transitions using Algorithm~\ref{alg:sumtree_sampling}
            \STATE Compute Importance Sampling weights $\tilde{w}_i$ via Equation~\ref{eq:per_is_weights}
            \FOR{each sample $i \in \{1, \dots, B\}$}
                \STATE $a^* \gets \arg\max_{a'} Q(s'_{i}, a'; \theta)$ \COMMENT{Action selection via online network}
                \STATE $Y_i \gets r_i + \gamma Q(s'_{i}, a^*; \theta^-)$ \COMMENT{Action evaluation via target network}
                \STATE $\delta_i \gets Q(s_i, a_i; \theta) - Y_i$ \COMMENT{Compute TD error}
                \STATE Update transition priority in SumTree: $p_i \gets (|\delta_i| + 10^{-5})^\alpha$
            \ENDFOR
            \STATE Compute weighted loss: $\mathcal{L} = \frac{1}{B}\sum_{i=1}^B \tilde{w}_i \cdot \mathcal{L}_{\text{Huber}}(\delta_i)$
            \STATE Execute Adam backpropagation: $\theta \gets \theta - \alpha \nabla_\theta \mathcal{L}$
            \STATE Execute Polyak soft update: $\theta^- \gets \tau \theta + (1 - \tau) \theta^-$
        \ENDIF
        \STATE $s_t \gets s_{t+1}$
    \ENDFOR
\ENDFOR
\RETURN $\theta^* \gets \theta$
\end{algorithmic}
\end{algorithm}
